%!TEX program = xelatex
% chktex-file 2
% chktex-file 12
% chktex-file 13
% chktex-file 38
% chktex-file 1
% chktex-file 8
% chktex-file 24
\documentclass[12pt,a4paper,oneside]{extreport}

\usepackage{fontspec}
\setmainfont{TeX Gyre Termes}
\setsansfont{TeX Gyre Termes}
\setmonofont{TeX Gyre Termes}
\usepackage{unicode-math}
\setmathfont{Latin Modern Math}

\usepackage[a4paper,top=2.5cm,bottom=3cm,left=3cm,right=2cm]{geometry}
\usepackage{setspace}
\setstretch{1.3}
\usepackage{indentfirst}
\setlength{\parindent}{1cm}
\setlength{\parskip}{6pt}

\usepackage{graphicx}
\usepackage[table]{xcolor}
\usepackage{array}
\usepackage{enumitem}
\usepackage{longtable}
\usepackage{booktabs}
\usepackage{tabularx}
\usepackage{multirow}
\usepackage{amsmath}
\usepackage{listings}
\usepackage{caption}
\usepackage{float}
\usepackage{tikz}
\usetikzlibrary{arrows.meta,calc,fit,positioning,shapes.geometric}
\usepackage{microtype}
\usepackage{titlesec}
\usepackage{fancyhdr}
\usepackage[hidelinks]{hyperref}
\usepackage{tocloft}
\urlstyle{same}
\emergencystretch=3em
\hfuzz=20pt
\hbadness=10000

\definecolor{reportgray}{HTML}{F3F4F6}
\definecolor{reportline}{HTML}{B0B7C3}

\renewcommand{\thechapter}{\Roman{chapter}}
\renewcommand{\chaptername}{}

\titleformat{\chapter}[display]
  {\normalfont\bfseries\color{black}}
  {\thechapter}
  {0.5em}
  {\Huge}
\titleformat{name=\chapter,numberless}[display]
  {\normalfont\bfseries\color{black}}
  {}
  {0pt}
  {\Huge}
\titleformat{\section}{\normalfont\LARGE\bfseries\color{black}}{\thesection}{0.8em}{}
\titleformat{\subsection}{\normalfont\Large\bfseries\color{black}}{\thesubsection}{0.8em}{}

\captionsetup{skip=6pt,font=small,labelfont=bf}
\setlength{\cftbeforetabskip}{4pt}
\setlength{\cftbeforefigskip}{4pt}
\cftsetindents{table}{0em}{3.2em}
\cftsetindents{figure}{0em}{3.2em}
\renewcommand{\cfttabfont}{\normalfont}
\renewcommand{\cfttabpagefont}{\normalfont}
\renewcommand{\cftfigfont}{\normalfont}
\renewcommand{\cftfigpagefont}{\normalfont}

\pagestyle{fancy}
\setlength{\headheight}{15pt}
\fancyhf{}
\fancyhead[L]{Leave Management System}
\fancyhead[R]{Báo cáo mini-project}
\fancyfoot[C]{\thepage}
\renewcommand{\headrulewidth}{0.4pt}

\lstdefinestyle{reportcode}{
  basicstyle=\small,
  breaklines=true,
  frame=single,
  rulecolor=\color{reportline},
  backgroundcolor=\color{reportgray},
  columns=fullflexible,
  keepspaces=true,
  showstringspaces=false
}
\lstset{style=reportcode}

\newcommand{\reporttitle}{Leave Management System}
\newcommand{\studentname}{Đinh Hiệp}
\newcommand{\studentemail}{hiepxuxi6688@gmail.com}
\newcommand{\mentorname}{Chưa cập nhật}
\newcommand{\unitname}{Viettel Digital Talent}

\begin{document}
\hypersetup{pageanchor=false}
\fontsize{13pt}{16.9pt}\selectfont

\begin{titlepage}
  \centering

  {\bfseries TẬP ĐOÀN CÔNG NGHIỆP - VIỄN THÔNG QUÂN ĐỘI\par}
  \vspace{1.2cm}

  {\Large \bfseries BÁO CÁO MINI-PROJECT\par}
  \vspace{1.5cm}

  {\Huge \bfseries \reporttitle\par}
  \vspace{1.5cm}

  {\Large \bfseries \studentname\par}
  \vspace{0.3cm}
  {\large \studentemail\par}

  \vfill

  {\large Chương trình Viettel Digital Talent 2026\par}
  \vspace{0.3cm}
  {\large Lĩnh vực: Software Engineer\par}

  \vspace{1.2cm}

  \begin{tabular}{rl}
    \textbf{Mentor:} & \mentorname \\
    \textbf{Đơn vị:} & \unitname \\
  \end{tabular}

  \vfill

  {\large HÀ NỘI - 2026\par}
\end{titlepage}

\pagenumbering{roman}
\hypersetup{pageanchor=true}

\phantomsection
\chapter*{Lời mở đầu}
\addcontentsline{toc}{chapter}{Lời mở đầu}
Báo cáo này trình bày quá trình phân tích, thiết kế, triển khai và đánh giá mini-project Leave Management System. Nội dung được xây dựng theo hướng báo cáo phát triển phần mềm web, tập trung vào bài toán quản lý nghỉ phép, kiến trúc backend/frontend, mô hình dữ liệu, luồng nghiệp vụ, kết quả giao diện, kiểm thử và các giới hạn của hệ thống trong phạm vi demo.

\phantomsection
\chapter*{Lời cảm ơn}
\addcontentsline{toc}{chapter}{Lời cảm ơn}
Em xin chân thành cảm ơn chương trình Viettel Digital Talent, mentor và các anh/chị đã hỗ trợ, góp ý trong quá trình thực hiện mini-project. Những phản hồi và định hướng trong quá trình phát triển giúp em hoàn thiện sản phẩm theo hướng rõ ràng, thực tế và có kỷ luật kỹ thuật hơn.

\phantomsection
\chapter*{Tóm tắt nội dung và đóng góp}
\addcontentsline{toc}{chapter}{Tóm tắt nội dung và đóng góp}
Leave Management System là hệ thống quản lý nghỉ phép dùng cho bối cảnh doanh nghiệp, bao gồm xác thực người dùng, quản lý phòng ban, loại nghỉ, quỹ phép, quy trình nộp và duyệt đơn, lịch tổng hợp, dashboard, báo cáo và thông báo. Hệ thống được xây dựng bằng Spring Boot, React, PostgreSQL và Docker, có phân quyền theo vai trò, migration cơ sở dữ liệu, kiểm thử tự động và dữ liệu demo phục vụ trình bày.

Đóng góp chính của mini-project là mô phỏng tương đối đầy đủ vòng đời phát triển một ứng dụng web nghiệp vụ: phân tích yêu cầu, mô hình hóa use case, thiết kế kiến trúc, thiết kế dữ liệu, hiện thực backend/frontend, kiểm thử và triển khai demo. Bên cạnh các chức năng cốt lõi, hệ thống còn có một số cải thiện phục vụ demo local như tệp đính kèm cho đơn nghỉ, báo cáo tổng hợp nâng cao, lịch theo phòng ban và dashboard được tinh chỉnh theo vai trò người dùng.

\phantomsection
\chapter*{Danh mục thuật ngữ và từ viết tắt}
\addcontentsline{toc}{chapter}{Danh mục thuật ngữ và từ viết tắt}
\begin{longtable}{p{0.22\textwidth}p{0.64\textwidth}}
\toprule
\textbf{Thuật ngữ} & \textbf{Giải thích} \\
\midrule
\endfirsthead
\toprule
\textbf{Thuật ngữ} & \textbf{Giải thích} \\
\midrule
\endhead
API & Giao diện lập trình ứng dụng, dùng để frontend giao tiếp với backend. \\
CRUD & Nhóm thao tác tạo, đọc, cập nhật và xóa dữ liệu. \\
DTO & Đối tượng truyền dữ liệu qua API, tách biệt với entity lưu trữ trong database. \\
ERD & Sơ đồ quan hệ thực thể, mô tả bảng dữ liệu và quan hệ khóa ngoại. \\
JWT & Token dùng cho xác thực stateless giữa frontend và backend. \\
RBAC & Cơ chế phân quyền dựa trên vai trò người dùng. \\
SPA & Ứng dụng web một trang, giao diện được render chủ yếu ở phía trình duyệt. \\
Smoke test & Kiểm thử nhanh các luồng chính để xác nhận hệ thống vẫn hoạt động sau thay đổi. \\
\bottomrule
\end{longtable}

\clearpage
\tableofcontents
\clearpage
\begingroup
\setstretch{1.0}
\setlength{\parskip}{0pt}
\phantomsection
\addcontentsline{toc}{chapter}{Danh mục hình vẽ}
\listoffigures
\endgroup
\clearpage
\begingroup
\setstretch{1.0}
\setlength{\parskip}{0pt}
\phantomsection
\addcontentsline{toc}{chapter}{Danh mục bảng biểu}
\listoftables
\endgroup
\clearpage

\pagenumbering{arabic}

\chapter{Giới thiệu}

\section{Đặt vấn đề}
Trong môi trường doanh nghiệp, nghỉ phép là một nghiệp vụ xuất hiện thường xuyên và liên quan trực tiếp đến nhiều bên: nhân viên cần chủ động đăng ký thời gian nghỉ, quản lý cần nắm được tình hình nhân sự để phê duyệt phù hợp, bộ phận nhân sự cần theo dõi quỹ phép và lưu lại lịch sử xử lý, còn ban quản trị cần có dữ liệu tổng hợp để đánh giá tình hình vận hành. Nếu quy trình này được xử lý bằng trao đổi miệng, tin nhắn rời rạc, bảng tính thủ công hoặc nhiều biểu mẫu không liên thông, thông tin dễ bị phân tán và khó kiểm soát.

Một quy trình nghỉ phép thủ công thường gặp một số vấn đề điển hình. Thứ nhất, người quản lý khó quan sát đầy đủ các đơn đang chờ xử lý, các khoảng thời gian nhân sự vắng mặt và khả năng chồng lấn giữa các đơn. Thứ hai, nhân viên không có một nơi thống nhất để tra cứu trạng thái đơn, lịch sử phê duyệt và số ngày phép còn lại. Thứ ba, bộ phận nhân sự phải tổng hợp dữ liệu từ nhiều nguồn, dễ sai lệch khi cập nhật quỹ phép, ngày lễ hoặc báo cáo theo tháng/quý. Thứ tư, khi thiếu lịch sử thao tác rõ ràng, việc truy vết ai đã phê duyệt, từ chối, hủy đơn hoặc điều chỉnh quỹ phép trở nên khó khăn.

Từ các vấn đề trên, một hệ thống quản lý nghỉ phép tập trung là cần thiết để số hóa toàn bộ vòng đời của đơn nghỉ: từ lúc nhân viên tạo đơn, hệ thống tự tính số ngày nghỉ hợp lệ, chuyển đơn đến người có thẩm quyền, ghi nhận quyết định phê duyệt hoặc từ chối, cập nhật quỹ phép và hiển thị thông tin trên lịch tổng hợp. Hệ thống cũng cần cung cấp dashboard, báo cáo và cơ chế thông báo để người dùng nắm được các sự kiện quan trọng mà không phải theo dõi thủ công.

Leave Management System được xây dựng như một mini-project phát triển phần mềm web nhằm mô phỏng bài toán trên trong phạm vi có kiểm soát. Dự án không chỉ tập trung vào việc tạo ra các màn hình thao tác, mà còn nhấn mạnh các khía cạnh cốt lõi của một ứng dụng web hoàn chỉnh: phân tích yêu cầu, thiết kế dữ liệu, thiết kế kiến trúc backend/frontend, quản lý phân quyền, kiểm thử, đóng gói bằng Docker và triển khai demo.

\section{Mục tiêu mini-project}
Mục tiêu tổng quát của mini-project là xây dựng một hệ thống web hỗ trợ quản lý quy trình nghỉ phép trong doanh nghiệp, có khả năng thể hiện đầy đủ các bước chính của một dự án phát triển phần mềm từ phân tích, thiết kế đến triển khai và đánh giá. Hệ thống cần có giao diện sử dụng được, API backend rõ ràng, cơ sở dữ liệu có ràng buộc nghiệp vụ, quy trình kiểm thử tự động và môi trường chạy nhất quán.

Về mặt nghiệp vụ, hệ thống hướng đến việc hỗ trợ nhân viên tạo, theo dõi, sửa hoặc hủy đơn nghỉ; hỗ trợ quản lý trực tiếp phê duyệt hoặc từ chối đơn; hỗ trợ bộ phận nhân sự theo dõi dữ liệu nghỉ phép toàn tổ chức; và hỗ trợ quản trị viên cấu hình các danh mục cần thiết như người dùng, phòng ban, loại nghỉ và quỹ phép. Ngoài ra, hệ thống cần cung cấp lịch tổng hợp, dashboard, báo cáo và thông báo để giúp thông tin được cập nhật kịp thời giữa các nhóm người dùng.

Về mặt kỹ thuật, dự án đặt mục tiêu xây dựng một ứng dụng web theo kiến trúc rõ ràng, tách biệt frontend, backend và database. Backend đảm nhiệm xử lý nghiệp vụ, bảo mật, phân quyền, kiểm tra dữ liệu và lưu trữ. Frontend cung cấp giao diện tương tác cho từng vai trò người dùng. Cơ sở dữ liệu lưu trữ các thực thể chính của bài toán và được quản lý thay đổi bằng migration. Toàn bộ hệ thống được container hóa để giảm khác biệt giữa môi trường phát triển và môi trường chạy thử.

Về mặt học thuật và báo cáo, mini-project được sử dụng để minh họa quy trình phát triển một phần mềm web có cấu trúc. Vì vậy, báo cáo không chỉ trình bày kết quả giao diện hoặc danh sách chức năng, mà còn phân tích cách xác định yêu cầu, mô hình hóa use case, thiết kế dữ liệu, mô tả luồng nghiệp vụ, lựa chọn kiến trúc, tổ chức mã nguồn và kiểm thử hệ thống.

\section{Phạm vi hệ thống}
Trong phạm vi hiện tại, Leave Management System tập trung vào các nhóm chức năng chính của quy trình quản lý nghỉ phép. Nhóm chức năng xác thực và phân quyền cho phép người dùng đăng nhập, sử dụng hệ thống theo vai trò và bảo vệ các thao tác nhạy cảm. Nhóm chức năng quản trị dữ liệu nền bao gồm quản lý người dùng, phòng ban, loại nghỉ phép, ngày lễ và quỹ phép. Đây là các dữ liệu làm cơ sở cho việc tạo đơn, tính số ngày nghỉ và xác định người có quyền xử lý.

Nhóm chức năng nghiệp vụ trung tâm là quản lý đơn nghỉ phép. Người dùng có thể tạo đơn với loại nghỉ, thời gian, phần ngày nghỉ và lý do; hệ thống tự tính số ngày nghỉ dựa trên ngày làm việc, cuối tuần, ngày lễ và lựa chọn nửa ngày. Đơn nghỉ sau đó đi qua quy trình xử lý gồm chờ duyệt, phê duyệt, từ chối hoặc hủy. Khi một đơn được xử lý, hệ thống ghi nhận lịch sử thao tác, cập nhật dữ liệu liên quan và bảo đảm các ràng buộc như không chồng lấn đơn hoặc không vượt quá quỹ phép còn lại.

Bên cạnh quy trình đơn nghỉ, hệ thống còn có các chức năng hỗ trợ quan sát và tổng hợp dữ liệu. Lịch nghỉ phép giúp người dùng nhìn được tình hình nghỉ theo tháng hoặc tuần. Dashboard cung cấp các chỉ số phù hợp với từng vai trò, chẳng hạn đơn đang chờ xử lý, số người đang nghỉ hoặc tình hình quỹ phép. Chức năng báo cáo hỗ trợ xuất dữ liệu phục vụ tổng hợp và phân tích. Cơ chế thông báo giúp người dùng nhận biết các thay đổi quan trọng trong quy trình xử lý đơn.

Một số hướng cải thiện có thể được xem xét sau khi hoàn thiện phần lõi, ví dụ duy trì tệp đính kèm cho đơn nghỉ ở môi trường demo local, bổ sung chế độ xem lịch theo phòng ban và tinh chỉnh trải nghiệm giao diện. Các hướng này không làm thay đổi mục tiêu cốt lõi của báo cáo: trình bày một hệ thống web quản lý nghỉ phép có quy trình nghiệp vụ rõ ràng, kiến trúc triển khai được và có khả năng mở rộng theo các nguyên tắc thiết kế đã chọn.

\section{Đối tượng sử dụng}
Hệ thống được thiết kế cho bốn nhóm người dùng chính. Nhân viên là nhóm sử dụng hệ thống ở góc nhìn self-service. Họ cần tạo đơn nghỉ, theo dõi trạng thái xử lý, xem lịch sử đơn của mình, tra cứu số ngày phép còn lại và quan sát lịch nghỉ liên quan đến nhóm hoặc phòng ban.

Quản lý trực tiếp là nhóm chịu trách nhiệm xử lý các đơn nghỉ của nhân viên thuộc phạm vi quản lý. Ngoài các chức năng của nhân viên, quản lý cần có danh sách đơn đang chờ duyệt, thông tin đủ rõ để đánh giá thời gian nghỉ, lý do nghỉ và khả năng ảnh hưởng đến nhân sự trong nhóm. Các quyết định phê duyệt hoặc từ chối của quản lý cần được ghi nhận lại để phục vụ truy vết.

Bộ phận nhân sự có góc nhìn tổng thể hơn đối với dữ liệu nghỉ phép. Nhóm này cần theo dõi lịch nghỉ của toàn tổ chức, quản lý một số dữ liệu nền như ngày lễ, điều chỉnh quỹ phép khi cần và khai thác báo cáo. Trong một số trường hợp, HR có thể cần can thiệp vào quy trình xử lý để đảm bảo dữ liệu hệ thống phản ánh đúng thực tế vận hành.

Quản trị viên chịu trách nhiệm cấu hình hệ thống và quản lý các dữ liệu có ảnh hưởng rộng, chẳng hạn tài khoản người dùng, vai trò, phòng ban, loại nghỉ và các thiết lập liên quan. Vai trò này giúp hệ thống có thể được khởi tạo, duy trì và điều chỉnh khi cơ cấu tổ chức hoặc chính sách nghỉ phép thay đổi. Chi tiết quyền hạn và các use case chính của từng nhóm người dùng sẽ được trình bày ở Chương II.

\section{Cấu trúc báo cáo}
Báo cáo được tổ chức theo mạch phát triển phần mềm từ bài toán đến giải pháp. Chương I giới thiệu bối cảnh, mục tiêu, phạm vi và các nhóm người dùng của hệ thống. Đây là phần đặt nền cho việc hiểu vì sao hệ thống cần được xây dựng và hệ thống hướng đến giải quyết những vấn đề nào.

Chương II trình bày nội dung phân tích và phương pháp xây dựng hệ thống. Phần này bao gồm phân tích yêu cầu, mô hình use case, kiến trúc tổng thể, thiết kế dữ liệu, các luồng nghiệp vụ quan trọng và phương pháp triển khai. Các sơ đồ như use case diagram, database diagram, state machine và sequence diagram được đặt trong chương này để thể hiện rõ cách hệ thống được mô hình hóa trước khi hiện thực.

Chương III trình bày kết quả thực hiện và đánh giá. Nội dung chương tập trung vào các chức năng đã hoàn thành, giao diện người dùng, kết quả triển khai, chiến lược kiểm thử và đánh giá mức độ đáp ứng mục tiêu ban đầu. Đây là phần liên kết giữa thiết kế ở Chương II và sản phẩm đã được xây dựng.

Chương IV tổng kết kết quả đạt được, nêu các bài học kỹ thuật, chỉ ra hạn chế của hệ thống và đề xuất hướng phát triển tiếp theo. Các nội dung chi tiết có tính chất tra cứu như danh sách API, schema cơ sở dữ liệu chi tiết hoặc hình ảnh giao diện có thể được đưa vào phụ lục để giữ phần thân báo cáo tập trung và mạch lạc.

\chapter{Nội dung và phương pháp}

\section{Phân tích yêu cầu}
Phân tích yêu cầu là bước xác định hệ thống cần hỗ trợ những nhóm nghiệp vụ nào, phục vụ các nhóm người dùng nào và phải đáp ứng những ràng buộc kỹ thuật nào để có thể vận hành ổn định. Đối với Leave Management System, yêu cầu được xây dựng xoay quanh quy trình nghỉ phép trong doanh nghiệp: nhân viên tạo đơn, người quản lý xử lý đơn, bộ phận nhân sự theo dõi dữ liệu toàn tổ chức và quản trị viên cấu hình các dữ liệu nền của hệ thống.

Trong phạm vi báo cáo, các yêu cầu được trình bày ở mức nhóm năng lực chính của hệ thống thay vì liệt kê toàn bộ chi tiết giao diện hoặc endpoint. Cách tiếp cận này giúp phần phân tích giữ được tính ổn định khi hệ thống tiếp tục mở rộng, đồng thời vẫn đủ cơ sở cho các phần thiết kế use case, kiến trúc, cơ sở dữ liệu và luồng nghiệp vụ ở các mục tiếp theo.

\subsection{Yêu cầu chức năng}
Yêu cầu chức năng mô tả các hành vi chính mà hệ thống cần cung cấp cho người dùng. Các chức năng này được nhóm theo năng lực nghiệp vụ để tránh phụ thuộc vào cách tổ chức màn hình hoặc API cụ thể.

\begin{longtable}{p{0.24\textwidth}p{0.66\textwidth}}
\toprule
\textbf{Nhóm chức năng} & \textbf{Mô tả yêu cầu} \\
\midrule
\endfirsthead
\toprule
\textbf{Nhóm chức năng} & \textbf{Mô tả yêu cầu} \\
\midrule
\endhead
Xác thực và phân quyền & Hệ thống cần cho phép người dùng đăng nhập, duy trì phiên làm việc an toàn và giới hạn thao tác theo vai trò. Các quyền truy cập phải phân biệt giữa nhân viên, quản lý, nhân sự và quản trị viên. \\
Quản lý dữ liệu tổ chức & Hệ thống cần hỗ trợ quản lý người dùng, phòng ban, vai trò, quan hệ quản lý trực tiếp và trạng thái hoạt động của tài khoản. Đây là nền tảng để xác định phạm vi xem dữ liệu và người có thẩm quyền phê duyệt. \\
Quản lý chính sách nghỉ phép & Hệ thống cần quản lý loại nghỉ, ngày lễ và quỹ phép của từng nhân viên theo từng loại nghỉ. Các thông tin này được dùng khi tính số ngày nghỉ hợp lệ và kiểm tra số ngày phép còn lại. \\
Quản lý đơn nghỉ & Người dùng cần tạo, xem, sửa hoặc hủy đơn nghỉ trong các điều kiện hợp lệ. Hệ thống cần kiểm tra ngày bắt đầu, ngày kết thúc, phần ngày nghỉ, lý do, sự chồng lấn với đơn khác và số ngày nghỉ được tính tự động. \\
Phê duyệt đơn nghỉ & Người quản lý hoặc người có thẩm quyền cần xem danh sách đơn chờ xử lý, phê duyệt hoặc từ chối đơn kèm ghi chú khi cần. Mỗi thao tác xử lý cần làm thay đổi trạng thái đơn theo đúng quy tắc nghiệp vụ. \\
Lịch nghỉ phép & Hệ thống cần hiển thị lịch nghỉ tổng hợp theo khoảng thời gian, hỗ trợ quan sát tình hình nghỉ của cá nhân, nhóm hoặc tổ chức tùy theo quyền truy cập. \\
Dashboard và báo cáo & Hệ thống cần cung cấp thông tin tổng hợp như số đơn đang chờ xử lý, tình hình nghỉ trong khoảng thời gian, quỹ phép và dữ liệu báo cáo để hỗ trợ theo dõi vận hành. \\
Thông báo & Hệ thống cần thông báo cho người dùng khi có sự kiện quan trọng liên quan đến đơn nghỉ, chẳng hạn đơn mới được tạo, được phê duyệt, bị từ chối hoặc bị hủy. \\
Quản trị hệ thống & Quản trị viên cần có khả năng cấu hình và duy trì các dữ liệu nền để hệ thống thích ứng với thay đổi của tổ chức hoặc chính sách nghỉ phép. \\
\bottomrule
\end{longtable}

Các nhóm chức năng trên tạo thành phạm vi nghiệp vụ cốt lõi của hệ thống. Những cải thiện như tệp đính kèm cho đơn nghỉ trong môi trường demo local, báo cáo nâng cao, thông báo rõ hơn, lịch theo phòng ban và tối ưu UI/UX có thể được bổ sung sau dựa trên cùng mô hình actor, dữ liệu và luồng xử lý.

\subsection{Yêu cầu phi chức năng}
Yêu cầu phi chức năng mô tả các tiêu chí chất lượng của hệ thống. Đây là phần quan trọng đối với một dự án phát triển phần mềm web, vì hệ thống không chỉ cần đúng nghiệp vụ mà còn phải có khả năng bảo trì, kiểm thử và triển khai nhất quán.

\begin{longtable}{p{0.24\textwidth}p{0.66\textwidth}}
\toprule
\textbf{Tiêu chí} & \textbf{Yêu cầu} \\
\midrule
\endfirsthead
\toprule
\textbf{Tiêu chí} & \textbf{Yêu cầu} \\
\midrule
\endhead
Bảo mật & Mật khẩu cần được lưu dưới dạng hash an toàn; API cần yêu cầu xác thực với các thao tác nghiệp vụ; quyền truy cập phải được kiểm tra ở cả cấp route và cấp nghiệp vụ quan trọng. \\
Phân quyền & Hệ thống cần áp dụng RBAC theo vai trò, đồng thời xét thêm phạm vi dữ liệu như đơn của chính mình, đơn của cấp dưới hoặc dữ liệu toàn tổ chức. \\
Audit và truy vết & Các thao tác quan trọng như phê duyệt, từ chối, hủy đơn hoặc điều chỉnh quỹ phép cần được ghi nhận để phục vụ kiểm tra lại sau này. \\
Tính đúng đắn nghiệp vụ & Quy tắc tính ngày nghỉ cần nhất quán giữa backend và giao diện, có xét cuối tuần, ngày lễ, nửa ngày và số ngày phép còn lại. \\
Khả năng bảo trì & Mã nguồn cần được tổ chức theo module chức năng, tách DTO khỏi entity, tách logic nghiệp vụ khỏi controller và quản lý schema bằng migration. \\
Khả năng kiểm thử & Các thành phần quan trọng cần có khả năng kiểm thử tự động ở nhiều mức: unit test cho logic nghiệp vụ, test tích hợp cho API/backend và smoke test cho luồng người dùng chính. \\
Khả năng triển khai & Hệ thống cần chạy được bằng môi trường container hóa, giảm phụ thuộc vào cấu hình máy cá nhân và hỗ trợ tách biệt cấu hình phát triển với cấu hình triển khai. \\
Khả năng sử dụng & Giao diện cần rõ ràng, phù hợp với vai trò người dùng, hỗ trợ trạng thái tải/rỗng/lỗi và hiển thị thông tin nghiệp vụ bằng tiếng Việt. \\
Múi giờ và dữ liệu ngày & Các tính toán ngày nghỉ, ngày lễ và dashboard cần thống nhất theo múi giờ Việt Nam để tránh sai lệch khi xử lý dữ liệu ngày. \\
\bottomrule
\end{longtable}

Các yêu cầu phi chức năng này định hướng các quyết định kỹ thuật như sử dụng JWT cho xác thực, BCrypt cho mật khẩu, Flyway cho migration, Docker cho môi trường chạy, và các tầng kiểm thử để giảm rủi ro khi thay đổi hệ thống.

\subsection{Phạm vi và giới hạn}
Phạm vi của mini-project tập trung vào việc mô phỏng một hệ thống quản lý nghỉ phép có đủ các thành phần chính của một ứng dụng web thực tế: giao diện người dùng, API backend, cơ sở dữ liệu quan hệ, xác thực, phân quyền, kiểm thử và môi trường triển khai. Hệ thống được thiết kế để phù hợp với bài toán doanh nghiệp quy mô vừa và nhỏ, nơi quy trình nghỉ phép vẫn có thể mô hình hóa bằng các vai trò và luồng phê duyệt tương đối rõ ràng.

Báo cáo tập trung vào các chức năng đã hình thành rõ trong hệ thống và các nguyên tắc thiết kế giúp hệ thống có thể tiếp tục mở rộng. Vì vậy, những nội dung như danh sách API đầy đủ, chi tiết mọi cột trong cơ sở dữ liệu hoặc ảnh chụp toàn bộ giao diện sẽ được ưu tiên đưa vào phụ lục hoặc các phần kết quả sau, thay vì làm nặng phần phân tích yêu cầu.

Một số hướng cải thiện có thể xuất hiện trong quá trình phát triển tiếp theo, chẳng hạn hỗ trợ tệp đính kèm cho đơn nghỉ ở môi trường demo local, bổ sung báo cáo nâng cao, cải thiện thông báo, hoàn thiện lịch theo phòng ban hoặc tối ưu UI/UX. Các hướng này không phải trọng tâm của phần phân tích yêu cầu hiện tại, nhưng có thể được mô tả như phần mở rộng tự nhiên của các nhóm chức năng đã nêu: quản lý đơn nghỉ, quản lý dữ liệu liên quan, phân quyền truy cập, tổng hợp dữ liệu và lưu vết thao tác.

Ngược lại, một số bài toán có phạm vi lớn hơn như multi-tenant, SSO doanh nghiệp hoặc đồng bộ hai chiều với lịch bên ngoài đòi hỏi thay đổi kiến trúc và giả định vận hành đáng kể. Các nội dung này không được xem là trọng tâm của mini-project, nhưng có thể được nhắc đến ở phần đánh giá hạn chế nếu cần làm rõ ranh giới giữa hệ thống minh họa và một sản phẩm doanh nghiệp quy mô lớn.

\section{Mô hình use case}
Mô hình use case mô tả cách các nhóm người dùng tương tác với hệ thống ở mức nghiệp vụ. Trong hệ thống quản lý nghỉ phép, bốn actor chính được xác định gồm Employee, Manager, HR và Admin. Employee là người sử dụng hệ thống để quản lý đơn nghỉ của bản thân. Manager xử lý đơn của cấp dưới và theo dõi lịch nghỉ của nhóm. HR theo dõi dữ liệu nghỉ phép toàn tổ chức, quản lý các chính sách liên quan đến nghỉ phép và khai thác báo cáo. Admin chịu trách nhiệm cấu hình dữ liệu nền và tài khoản hệ thống.

\subsection{Use case diagram}
Sơ đồ use case được giữ ở mức tổng quan để thể hiện phạm vi chính của hệ thống. Các use case nhỏ như lọc danh sách, xem chi tiết hoặc đánh dấu thông báo đã đọc không được tách riêng trên sơ đồ nhằm tránh làm rối mô hình.

\begin{figure}[H]
\centering
\resizebox{\textwidth}{!}{%
\begin{tikzpicture}[
  font=\small,
  card/.style={draw, rounded corners, fill=reportgray, text width=3.15cm, minimum height=4.1cm, align=left, inner sep=7pt},
  title/.style={font=\small\bfseries, align=center}
]
  \node[title] at (5.9,4.9) {Hệ thống quản lý nghỉ phép};

  \node[card] (employee) at (0,2.25) {
    \centering\textbf{Employee}\par
    \raggedright
    \vspace{0.15cm}
    Đăng nhập\\
    Quản lý hồ sơ\\
    Nộp/sửa/hủy đơn\\
    Xem lịch nghỉ cá nhân\\
    Nhận thông báo
  };

  \node[card] (manager) at (3.95,2.25) {
    \centering\textbf{Manager}\par
    \raggedright
    \vspace{0.15cm}
    Duyệt/từ chối đơn\\
    Theo dõi lịch nhóm\\
    Xem dashboard nhóm\\
    Nhận thông báo xử lý
  };

  \node[card] (hr) at (7.9,2.25) {
    \centering\textbf{HR}\par
    \raggedright
    \vspace{0.15cm}
    Theo dõi toàn tổ chức\\
    Quản lý chính sách nghỉ\\
    Quản lý quỹ phép\\
    Xuất báo cáo
  };

  \node[card] (admin) at (11.85,2.25) {
    \centering\textbf{Admin}\par
    \raggedright
    \vspace{0.15cm}
    Quản lý người dùng\\
    Quản lý phòng ban\\
    Quản lý danh mục nền\\
    Cấu hình hệ thống
  };

  \node[draw, dashed, rounded corners, fit=(employee)(manager)(hr)(admin), inner sep=0.35cm] {};
\end{tikzpicture}%
}
\caption{Use case diagram tổng quan}
\end{figure}

\subsection{Danh sách use case chính}
Bảng dưới đây tóm tắt các use case chính của hệ thống. Mỗi use case được mô tả ở mức mục tiêu nghiệp vụ, actor chính và kết quả mong đợi.

\begin{longtable}{p{0.20\textwidth}p{0.19\textwidth}p{0.47\textwidth}}
\toprule
\textbf{Use case} & \textbf{Actor chính} & \textbf{Mục tiêu và kết quả} \\
\midrule
\endfirsthead
\toprule
\textbf{Use case} & \textbf{Actor chính} & \textbf{Mục tiêu và kết quả} \\
\midrule
\endhead
Đăng nhập & Tất cả actor & Người dùng xác thực bằng email và mật khẩu. Hệ thống trả về phiên làm việc hợp lệ, thông tin vai trò và quyền truy cập tương ứng. \\
Quản lý hồ sơ cá nhân & Employee, Manager, HR, Admin & Người dùng xem thông tin cá nhân, phòng ban, vai trò, quản lý trực tiếp và các dữ liệu liên quan đến tài khoản của mình. \\
Nộp đơn nghỉ & Employee & Nhân viên nhập loại nghỉ, khoảng thời gian, phần ngày và lý do. Hệ thống kiểm tra dữ liệu, tính số ngày nghỉ, kiểm tra chồng lấn và tạo đơn ở trạng thái chờ duyệt. \\
Sửa hoặc hủy đơn & Employee & Nhân viên chỉnh sửa đơn còn chờ duyệt hoặc hủy đơn trong điều kiện hợp lệ. Hệ thống cập nhật trạng thái, ghi lịch sử thao tác và gửi thông báo cho bên liên quan. \\
Duyệt hoặc từ chối đơn & Manager, HR, Admin & Người có thẩm quyền xử lý đơn chờ duyệt. Khi duyệt, hệ thống cập nhật trạng thái và quỹ phép nếu loại nghỉ yêu cầu trừ phép; khi từ chối, hệ thống yêu cầu ghi chú lý do. \\
Xem lịch nghỉ & Employee, Manager, HR, Admin & Người dùng xem lịch nghỉ theo phạm vi quyền: cá nhân, nhóm hoặc toàn tổ chức. Hệ thống hiển thị các đơn đã duyệt và có thể kèm đơn đang chờ xử lý. \\
Quản lý danh mục & HR, Admin & Người dùng có quyền quản lý phòng ban, ngày lễ, loại nghỉ và các dữ liệu nền cần thiết cho nghiệp vụ nghỉ phép. \\
Quản lý quỹ phép & HR, Admin & Người dùng có quyền xem và điều chỉnh quỹ phép của nhân viên theo năm và loại nghỉ, phục vụ chính sách nghỉ phép của tổ chức. \\
Xem dashboard và báo cáo & Manager, HR, Admin & Người dùng theo dõi số liệu tổng hợp như đơn chờ duyệt, tình hình nghỉ theo thời gian, dữ liệu quỹ phép và xuất báo cáo khi cần. \\
Nhận thông báo & Employee, Manager, HR & Người dùng nhận thông báo khi có sự kiện liên quan đến đơn nghỉ, ví dụ đơn mới, đơn được cập nhật, được duyệt, bị từ chối hoặc bị hủy. \\
\bottomrule
\end{longtable}

\section{Thiết kế kiến trúc hệ thống}
Leave Management System được thiết kế theo mô hình web application ba phần chính: client chạy trên trình duyệt, backend cung cấp REST API và cơ sở dữ liệu PostgreSQL. Cách tổ chức này phù hợp với mục tiêu của dự án vì tách rõ trách nhiệm giữa giao diện, nghiệp vụ và lưu trữ dữ liệu, đồng thời vẫn giữ kiến trúc đủ đơn giản để phát triển và kiểm thử trong phạm vi mini-project.

\subsection{Architecture diagram}
Sơ đồ kiến trúc tổng quan thể hiện luồng giao tiếp chính và các concern cắt ngang của hệ thống.

\begin{figure}[H]
\centering
\resizebox{0.88\textwidth}{!}{%
\begin{tikzpicture}[
  font=\scriptsize,
  >=Latex,
  tier/.style={draw, rounded corners, fill=reportgray, text width=11.0cm, minimum height=1.0cm, align=center, inner sep=7pt},
  strip/.style={draw, rounded corners, fill=white, text width=9.8cm, minimum height=0.48cm, align=center, inner sep=3pt},
  arrow/.style={-Latex, line width=0.5pt}
]
  \node[tier] (frontend) at (0,7.0) {
    \textbf{Frontend}\\
    React SPA, TypeScript, Vite, Tailwind, TanStack Query, React Hook Form, Zod
  };

  \node[draw, rounded corners, fill=reportgray, minimum width=11.6cm, minimum height=3.2cm] (backend) at (0,4.0) {};
  \node[font=\small\bfseries] at (0,5.28) {Backend Spring Boot};
  \node[strip] at (0,4.75) {Controller layer -- REST API};
  \node[strip] at (0,4.25) {Service layer -- business rules and transactions};
  \node[strip] at (0,3.75) {Repository layer -- Spring Data JPA};
  \node[strip] at (0,3.25) {Domain layer -- Entity, DTO, Mapper};
  \node[strip] at (0,2.75) {Cross-cutting -- Security, Validation, Audit, Notification};

  \node[tier] (database) at (0,1.0) {
    \textbf{Database}\\
    PostgreSQL, Flyway migration, constraints, indexes, transactional data
  };

  \draw[arrow] (frontend.south) -- (backend.north);
  \draw[arrow] (backend.south) -- (database.north);
\end{tikzpicture}%
}
\caption{Architecture diagram tổng quan}
\end{figure}

\subsection{Backend architecture}
Backend được triển khai bằng Spring Boot theo hướng monolith có tổ chức module theo feature. Các package như \texttt{auth}, \texttt{user}, \texttt{department}, \texttt{leavetype}, \texttt{leavebalance}, \texttt{leaverequest}, \texttt{holiday}, \texttt{notification} và \texttt{report} phản ánh trực tiếp các vùng nghiệp vụ của hệ thống. Cách tổ chức này giúp mã nguồn dễ tìm, dễ kiểm thử và tránh việc một module nghiệp vụ bị phân tán quá nhiều theo lớp kỹ thuật.

Bên trong mỗi feature, hệ thống vẫn duy trì phân lớp quen thuộc gồm controller, service, repository, domain và DTO. Controller chịu trách nhiệm nhận request, validate dữ liệu đầu vào và trả response chuẩn. Service chứa logic nghiệp vụ như tính ngày nghỉ, kiểm tra chồng lấn, cập nhật trạng thái đơn và điều chỉnh quỹ phép. Repository dùng Spring Data JPA để truy vấn dữ liệu. Entity chỉ đại diện cho mô hình lưu trữ, không được trả trực tiếp qua API; dữ liệu ra vào hệ thống đi qua DTO để giảm rò rỉ chi tiết nội bộ.

Các concern dùng chung như response wrapper, exception handler, audit log, cấu hình bảo mật và auditing được đặt trong các package chung hoặc cấu hình riêng. Nhờ vậy, luồng xử lý nghiệp vụ chính vẫn tập trung trong service, còn các quy tắc kỹ thuật lặp lại được xử lý nhất quán ở một nơi.

\subsection{Frontend architecture}
Frontend được triển khai bằng React, TypeScript và Vite. Cấu trúc mã nguồn đi theo feature, tương ứng với các vùng chức năng như đăng nhập, người dùng, phòng ban, loại nghỉ, quỹ phép, đơn nghỉ, phê duyệt, lịch, dashboard, báo cáo và thông báo. Mỗi feature chứa các thành phần giao diện, hook, API client và kiểu dữ liệu liên quan.

Server state được quản lý bằng TanStack Query để chuẩn hóa việc gọi API, cache, refetch và xử lý trạng thái tải/lỗi. Form được xây dựng với React Hook Form và Zod để kiểm soát validation ở phía client trước khi gửi request. Cách này không thay thế validation ở backend, nhưng giúp giao diện phản hồi nhanh hơn và giảm request sai định dạng.

Routing phía frontend được bảo vệ theo trạng thái đăng nhập và vai trò người dùng. Các màn hình quản trị, nhân sự hoặc phê duyệt chỉ được hiển thị cho actor phù hợp. Việc phân quyền cuối cùng vẫn được backend kiểm tra, nhưng frontend giúp điều hướng người dùng đúng phạm vi chức năng và giảm thao tác không hợp lệ.

\subsection{Security architecture}
Hệ thống sử dụng JWT cho xác thực API theo mô hình stateless. Sau khi đăng nhập thành công, client gửi access token trong header Authorization cho các request cần bảo vệ. Refresh token được lưu phía backend dưới dạng hash để có thể thu hồi khi người dùng đăng xuất hoặc token hết hạn.

Mật khẩu được mã hóa bằng BCrypt trước khi lưu trữ. Spring Security kiểm soát filter chain, CORS, xác thực request và phân quyền ở cấp endpoint. Bên cạnh RBAC theo vai trò, một số thao tác còn kiểm tra phạm vi dữ liệu ở tầng service, ví dụ nhân viên chỉ được sửa đơn của chính mình, manager chỉ xử lý đơn thuộc phạm vi quản lý, còn HR/Admin có phạm vi rộng hơn.

Các thao tác quan trọng như duyệt, từ chối, hủy đơn và điều chỉnh dữ liệu nghiệp vụ được ghi lại thông qua approval action hoặc audit log. Điều này giúp hệ thống có khả năng truy vết nguyên nhân thay đổi trạng thái, không chỉ lưu trạng thái cuối cùng.

\section{Thiết kế dữ liệu}
Mô hình dữ liệu được thiết kế trên PostgreSQL và được quản lý bằng Flyway migration. Cơ sở dữ liệu tập trung vào ba nhóm thông tin chính: dữ liệu tổ chức, dữ liệu chính sách nghỉ phép và dữ liệu phát sinh từ quy trình đơn nghỉ. Ngoài ra, hệ thống có các bảng hỗ trợ cho xác thực, thông báo và audit.

\subsection{Database diagram}
Sơ đồ ERD dưới đây chỉ thể hiện các bảng chính và quan hệ khóa ngoại quan trọng. Các cột audit dùng chung như \texttt{created\_at}, \texttt{updated\_at}, \texttt{created\_by}, \texttt{updated\_by} không được lặp lại trong sơ đồ để giữ mô hình dễ đọc.

\begin{figure}[H]
\centering
\resizebox{\textwidth}{!}{%
\begin{tikzpicture}[
  font=\scriptsize,
  >=Latex,
  entity/.style={draw, rounded corners, fill=reportgray, text width=3.35cm, minimum height=1.65cm, align=left, inner sep=5pt},
  rel/.style={-Latex, line width=0.42pt},
  softrel/.style={-Latex, dashed, line width=0.38pt},
  note/.style={font=\scriptsize, align=center}
]
  \node[entity] (departments) at (-5.0,3.0) {
    \centering\textbf{\texttt{departments}}\par
    \raggedright
    \texttt{id}\\
    \texttt{head\_user\_id -> users}
  };

  \node[entity] (users) at (0,3.0) {
    \centering\textbf{\texttt{users}}\par
    \raggedright
    \texttt{id}\\
    \texttt{department\_id -> departments}\\
    \texttt{manager\_id -> users}
  };

  \node[entity] (tokens) at (5.0,3.0) {
    \centering\textbf{\texttt{refresh\_tokens}}\par
    \raggedright
    \texttt{id}\\
    \texttt{user\_id -> users}
  };

  \node[entity] (types) at (-5.0,0.55) {
    \centering\textbf{\texttt{leave\_types}}\par
    \raggedright
    \texttt{id}\\
    \texttt{default\_quota}\\
    \texttt{requires\_balance}
  };

  \node[entity] (requests) at (0,0.55) {
    \centering\textbf{\texttt{leave\_requests}}\par
    \raggedright
    \texttt{user\_id -> users}\\
    \texttt{leave\_type\_id -> leave\_types}\\
    \texttt{manager\_id -> users}
  };

  \node[entity] (actions) at (5.0,0.55) {
    \centering\textbf{\texttt{approval\_actions}}\par
    \raggedright
    \texttt{leave\_request\_id -> leave\_requests}\\
    \texttt{actor\_id -> users}
  };

  \node[entity] (balances) at (-5.0,-1.9) {
    \centering\textbf{\texttt{leave\_balances}}\par
    \raggedright
    \texttt{user\_id -> users}\\
    \texttt{leave\_type\_id -> leave\_types}\\
    \texttt{year}
  };

  \node[entity] (notifications) at (0,-1.9) {
    \centering\textbf{\texttt{notifications}}\par
    \raggedright
    \texttt{user\_id -> users}\\
    \texttt{leave\_request\_id -> leave\_requests}
  };

  \node[entity] (audit) at (5.0,-1.9) {
    \centering\textbf{\texttt{audit\_log}}\par
    \raggedright
    \texttt{actor\_id -> users}\\
    \texttt{target\_type, target\_id}
  };

  \node[entity, fill=white] (holidays) at (-5.0,-4.05) {
    \centering\textbf{\texttt{holidays}}\par
    \raggedright
    Bảng tham chiếu ngày lễ, dùng khi tính số ngày nghỉ.
  };

  \draw[rel] (users.west) -- (departments.east);
  \draw[rel] (tokens.west) -- (users.east);
  \draw[rel] (requests.north) -- (users.south);
  \draw[rel] (requests.west) -- (types.east);
  \draw[rel] (actions.west) -- (requests.east);
  \draw[rel] (balances.north) -- (types.south);
  \draw[rel] (notifications.north) -- (requests.south);
  \draw[rel] (audit.north west) to[out=135,in=-35] (users.south east);
  \draw[softrel] (holidays.east) -- (requests.south west);
\end{tikzpicture}%
}
\caption{Database diagram rút gọn với khóa ngoại chính}
\end{figure}

\subsection{Các bảng dữ liệu chính}
Bảng sau mô tả vai trò của các bảng chính trong cơ sở dữ liệu. Báo cáo chỉ trình bày ở mức ý nghĩa nghiệp vụ; chi tiết từng cột có thể được đưa vào phụ lục nếu cần đối chiếu với migration.

\begin{longtable}{p{0.25\textwidth}p{0.62\textwidth}}
\toprule
\textbf{Bảng} & \textbf{Vai trò trong hệ thống} \\
\midrule
\endfirsthead
\toprule
\textbf{Bảng} & \textbf{Vai trò trong hệ thống} \\
\midrule
\endhead
\texttt{departments} & Lưu phòng ban, mã phòng ban, tên phòng ban và người phụ trách nếu có. Bảng này là cơ sở để nhóm người dùng và phân tích dữ liệu theo đơn vị tổ chức. \\
\texttt{users} & Lưu tài khoản người dùng, thông tin nhân viên, vai trò, phòng ban, quản lý trực tiếp và trạng thái hoạt động. Đây là bảng trung tâm cho xác thực và phân quyền. \\
\texttt{leave\_types} & Lưu các loại nghỉ, quota mặc định và cờ xác định loại nghỉ có cần kiểm tra quỹ phép hay không. \\
\texttt{leave\_balances} & Lưu quỹ phép của từng người dùng theo loại nghỉ và năm, bao gồm tổng ngày, số ngày đã dùng, số ngày điều chỉnh và số ngày chuyển tiếp nếu có. \\
\texttt{holidays} & Lưu ngày lễ để loại khỏi số ngày làm việc khi tính thời lượng nghỉ phép. \\
\texttt{leave\_requests} & Lưu đơn nghỉ, người tạo, loại nghỉ, khoảng thời gian, nửa ngày, tổng ngày nghỉ, lý do, trạng thái và người quản lý xử lý. \\
\texttt{approval\_actions} & Lưu lịch sử hành động trên đơn nghỉ như tạo, cập nhật, duyệt, từ chối hoặc hủy. Bảng này giúp truy vết diễn tiến của từng đơn. \\
\texttt{notifications} & Lưu thông báo trong ứng dụng cho các sự kiện vòng đời đơn nghỉ, gắn với người nhận và đơn liên quan nếu có. \\
\texttt{audit\_log} & Lưu vết thao tác quan trọng ở mức hệ thống, gồm actor, loại đối tượng, định danh đối tượng và dữ liệu trước/sau khi thay đổi. \\
\texttt{refresh\_tokens} & Lưu refresh token đã hash, thời điểm hết hạn và thời điểm thu hồi để hỗ trợ cơ chế đăng nhập stateless nhưng vẫn có khả năng revoke token. \\
\bottomrule
\end{longtable}

\subsection{Quản lý migration}
Schema cơ sở dữ liệu được quản lý bằng Flyway thay vì để JPA tự sinh bảng. Mỗi thay đổi schema được viết thành một migration SQL có thứ tự, giúp môi trường phát triển, kiểm thử và triển khai cùng đi qua một lịch sử thay đổi giống nhau. Cách tiếp cận này giảm rủi ro sai lệch schema giữa các môi trường và giúp dễ truy vết thời điểm một bảng, cột hoặc constraint được bổ sung.

Trong hệ thống, migration nền tạo các bảng cốt lõi như phòng ban, người dùng, loại nghỉ, quỹ phép, ngày lễ, đơn nghỉ, lịch sử phê duyệt, audit log và refresh token. Các migration sau bổ sung dữ liệu tham chiếu hoặc mở rộng schema cho những năng lực mới như thông báo và chuyển tiếp quỹ phép. Khi cần thay đổi cấu trúc dữ liệu, hệ thống tạo migration mới thay vì sửa migration đã áp dụng để tránh lỗi checksum và bảo đảm tính lặp lại của quá trình khởi tạo database.

\section{Thiết kế luồng nghiệp vụ}
Luồng nghiệp vụ trung tâm của hệ thống là vòng đời đơn nghỉ phép. Một đơn nghỉ được tạo ở trạng thái chờ duyệt, sau đó có thể được phê duyệt, từ chối hoặc hủy tùy theo actor và điều kiện nghiệp vụ. Việc thiết kế rõ trạng thái và luồng xử lý giúp backend kiểm soát tính hợp lệ của thao tác, đồng thời giúp frontend hiển thị đúng hành động người dùng có thể thực hiện.

\subsection{State machine của đơn nghỉ}
State machine của đơn nghỉ gồm bốn trạng thái chính: \texttt{PENDING}, \texttt{APPROVED}, \texttt{REJECTED} và \texttt{CANCELLED}. Trạng thái \texttt{PENDING} là điểm bắt đầu sau khi nhân viên nộp đơn. Từ trạng thái này, đơn có thể được chỉnh sửa, được duyệt, bị từ chối hoặc bị hủy. Khi đơn đã được duyệt, hệ thống cho phép hủy trong các điều kiện được kiểm soát và khôi phục quỹ phép nếu loại nghỉ có trừ quỹ.

\begin{figure}[H]
\centering
\resizebox{\textwidth}{!}{%
\begin{tikzpicture}[
  font=\scriptsize,
  >=Latex,
  state/.style={draw, rounded corners, fill=reportgray, minimum width=2.9cm, minimum height=0.95cm, align=center},
  label/.style={font=\scriptsize, fill=white, inner sep=1.5pt, align=center},
  arrow/.style={-Latex, line width=0.5pt}
]
  \node[state] (pending) at (0,0) {\texttt{PENDING}\\Chờ duyệt};
  \node[state] (approved) at (5.2,2.0) {\texttt{APPROVED}\\Đã duyệt};
  \node[state] (rejected) at (5.2,-2.0) {\texttt{REJECTED}\\Bị từ chối};
  \node[state] (cancelled) at (10.4,0) {\texttt{CANCELLED}\\Đã hủy};

  \draw[arrow] (-3.2,0) -- node[label, above]{Tạo đơn hợp lệ} (pending.west);
  \draw[arrow] (pending.north west) to[out=135,in=45,looseness=1.9]
    node[label, above]{Sửa khi còn chờ} (pending.north east);
  \draw[arrow] (pending.east) -- node[label, above, pos=0.46]{Duyệt} (approved.west);
  \draw[arrow] (pending.east) -- node[label, below, pos=0.46]{Từ chối} (rejected.west);
  \draw[arrow] (pending.south) .. controls (2.0,-3.25) and (8.2,-3.25) ..
    node[label, below, pos=0.5]{Hủy khi còn chờ} (cancelled.south);
  \draw[arrow] (approved.east) .. controls (7.1,2.75) and (9.3,2.2) ..
    node[label, above, pos=0.52]{Hủy có điều kiện\\và hoàn quỹ} (cancelled.north west);
\end{tikzpicture}%
}
\caption{State machine của đơn nghỉ}
\end{figure}

Các trạng thái \texttt{REJECTED} và \texttt{CANCELLED} là trạng thái kết thúc của một đơn nghỉ trong phạm vi hiện tại. Hệ thống không chuyển một đơn đã từ chối hoặc đã hủy quay lại trạng thái chờ duyệt; nếu người dùng vẫn có nhu cầu nghỉ, họ cần tạo đơn mới để hệ thống tính lại dữ liệu và kiểm tra điều kiện tại thời điểm mới.

\subsection{Sequence diagram: nộp đơn nghỉ}
Luồng nộp đơn nghỉ bắt đầu từ giao diện nhân viên. Backend không tin cậy hoàn toàn dữ liệu từ frontend mà thực hiện lại toàn bộ kiểm tra quan trọng, bao gồm khoảng ngày, nửa ngày, ngày lễ, chồng lấn đơn và quỹ phép.

\begin{figure}[H]
\centering
\resizebox{\textwidth}{!}{%
\begin{tikzpicture}[
  font=\scriptsize,
  >=Latex,
  msg/.style={-Latex, line width=0.45pt},
  life/.style={dashed, line width=0.35pt},
  actor/.style={draw, rounded corners, fill=reportgray, minimum width=2.2cm, minimum height=0.65cm, align=center}
]
  \node[actor] (fe) at (0,0) {Employee\\Frontend};
  \node[actor] (api) at (3.4,0) {LeaveRequest\\Controller};
  \node[actor] (svc) at (6.8,0) {LeaveRequest\\Service};
  \node[actor] (domain) at (10.2,0) {Calculator\\Repository};
  \node[actor] (notify) at (13.6,0) {Notification};

  \draw[life] (0,-0.4) -- (0,-6.6);
  \draw[life] (3.4,-0.4) -- (3.4,-6.6);
  \draw[life] (6.8,-0.4) -- (6.8,-6.6);
  \draw[life] (10.2,-0.4) -- (10.2,-6.6);
  \draw[life] (13.6,-0.4) -- (13.6,-6.6);

  \draw[msg] (0,-1.0) -- node[above]{Gửi đơn nghỉ} (3.4,-1.0);
  \draw[msg] (3.4,-1.7) -- node[above]{Validate DTO, lấy principal} (6.8,-1.7);
  \draw[msg] (6.8,-2.4) -- node[above]{Lấy user, loại nghỉ, ngày lễ} (10.2,-2.4);
  \draw[msg] (10.2,-3.1) -- node[above]{Tính số ngày nghỉ} (6.8,-3.1);
  \draw[msg] (6.8,-3.8) -- node[above]{Kiểm tra overlap và balance} (10.2,-3.8);
  \draw[msg] (6.8,-4.5) -- node[above]{Lưu \texttt{leave\_requests}} (10.2,-4.5);
  \draw[msg] (6.8,-5.2) -- node[above]{Ghi \texttt{approval\_actions}} (10.2,-5.2);
  \draw[msg] (6.8,-5.9) -- node[above]{Tạo thông báo cho manager} (13.6,-5.9);
  \draw[msg] (6.8,-6.5) -- node[above]{Trả đơn ở trạng thái \texttt{PENDING}} (0,-6.5);
\end{tikzpicture}%
}
\caption{Sequence diagram cho luồng nộp đơn nghỉ}
\end{figure}

Kết quả của luồng này là một bản ghi đơn nghỉ ở trạng thái \texttt{PENDING}, một bản ghi lịch sử hành động \texttt{CREATED} và thông báo gửi đến người quản lý trực tiếp. Nếu bất kỳ điều kiện nào không hợp lệ, service trả lỗi nghiệp vụ và không tạo đơn.

\subsection{Sequence diagram: duyệt hoặc từ chối đơn}
Luồng duyệt hoặc từ chối đơn được thực hiện bởi Manager, HR hoặc Admin tùy phạm vi quyền. Backend kiểm tra trạng thái hiện tại của đơn trước khi chuyển trạng thái để tránh xử lý lại đơn đã kết thúc.

\begin{figure}[H]
\centering
\resizebox{\textwidth}{!}{%
\begin{tikzpicture}[
  font=\scriptsize,
  >=Latex,
  msg/.style={-Latex, line width=0.45pt},
  life/.style={dashed, line width=0.35pt},
  actor/.style={draw, rounded corners, fill=reportgray, minimum width=2.25cm, minimum height=0.65cm, align=center}
]
  \node[actor] (mgr) at (0,0) {Manager/HR\\Frontend};
  \node[actor] (api) at (3.4,0) {Approval\\Endpoint};
  \node[actor] (svc) at (6.8,0) {LeaveRequest\\Service};
  \node[actor] (bal) at (10.2,0) {Balance\\Repository};
  \node[actor] (audit) at (13.6,0) {Audit\\Notification};

  \draw[life] (0,-0.4) -- (0,-6.8);
  \draw[life] (3.4,-0.4) -- (3.4,-6.8);
  \draw[life] (6.8,-0.4) -- (6.8,-6.8);
  \draw[life] (10.2,-0.4) -- (10.2,-6.8);
  \draw[life] (13.6,-0.4) -- (13.6,-6.8);

  \draw[msg] (0,-1.0) -- node[above]{Chọn duyệt hoặc từ chối} (3.4,-1.0);
  \draw[msg] (3.4,-1.7) -- node[above]{Kiểm tra quyền và principal} (6.8,-1.7);
  \draw[msg] (6.8,-2.4) to[loop right, looseness=5] node[right, align=left]{Tải đơn\\yêu cầu \texttt{PENDING}} (6.8,-2.4);
  \draw[msg] (6.8,-3.1) -- node[above]{Nếu duyệt: trừ quỹ phép} (10.2,-3.1);
  \draw[msg] (6.8,-3.8) -- node[above]{Cập nhật trạng thái} (10.2,-3.8);
  \draw[msg] (6.8,-4.5) -- node[above]{Ghi approval action} (10.2,-4.5);
  \draw[msg] (6.8,-5.2) -- node[above]{Ghi audit log} (13.6,-5.2);
  \draw[msg] (6.8,-5.9) -- node[above]{Thông báo kết quả cho requester} (13.6,-5.9);
  \draw[msg] (6.8,-6.6) -- node[above]{Trả trạng thái mới} (0,-6.6);
\end{tikzpicture}%
}
\caption{Sequence diagram cho luồng duyệt hoặc từ chối đơn}
\end{figure}

Trong trường hợp phê duyệt, hệ thống thực hiện kiểm tra quỹ phép ở thời điểm duyệt để bảo đảm dữ liệu vẫn hợp lệ sau khi có thể đã phát sinh các đơn khác. Trong trường hợp từ chối, hệ thống yêu cầu ghi chú lý do và không thay đổi quỹ phép. Cả hai nhánh đều ghi lịch sử hành động và tạo thông báo cho người nộp đơn.

\section{Phương pháp triển khai}
Phương pháp triển khai của dự án đi theo từng lát cắt chức năng, bắt đầu từ nền tảng hạ tầng và dữ liệu, sau đó phát triển dần các nhóm nghiệp vụ cốt lõi. Cách tiếp cận này giúp mỗi phần được kiểm chứng sớm bằng API, giao diện và test tương ứng, thay vì dồn toàn bộ kiểm thử về cuối.

\subsection{Công nghệ sử dụng}
Backend sử dụng Java 21 và Spring Boot 3.3. Spring Security đảm nhiệm xác thực và phân quyền, Spring Data JPA hỗ trợ truy cập dữ liệu, Flyway quản lý migration cơ sở dữ liệu và PostgreSQL là hệ quản trị cơ sở dữ liệu quan hệ chính. Các thành phần nghiệp vụ được tổ chức theo feature để giữ mã nguồn gần với domain của hệ thống.

Frontend sử dụng React 18, TypeScript, Vite, Tailwind CSS và các component theo hướng tái sử dụng. TanStack Query được dùng cho server state, React Hook Form và Zod được dùng cho form validation. Các request từ frontend đi qua lớp API client thay vì gọi trực tiếp trong component, giúp giao diện dễ kiểm thử và dễ thay đổi khi backend API điều chỉnh.

Môi trường phát triển và chạy thử được container hóa bằng Docker Compose. Cách này giúp backend, frontend và database chạy trong cùng một cấu hình có thể lặp lại. CI được dùng để kiểm tra các bước chất lượng như build, test và lint trước khi hợp nhất thay đổi.

\subsection{Quy trình phát triển}
Quy trình phát triển được chia theo các nhóm năng lực tăng dần. Giai đoạn đầu tập trung dựng nền tảng gồm Docker Compose, cấu trúc backend/frontend, schema cơ sở dữ liệu và migration ban đầu. Sau đó hệ thống bổ sung xác thực, phân quyền, quản lý người dùng, phòng ban, loại nghỉ, quỹ phép và ngày lễ.

Khi nền tảng dữ liệu đã ổn định, trọng tâm chuyển sang nghiệp vụ đơn nghỉ: nộp đơn, tính ngày nghỉ, kiểm tra ngày lễ/cuối tuần, kiểm tra chồng lấn, phê duyệt, từ chối, hủy đơn, ghi lịch sử hành động và audit. Các chức năng quan sát như lịch tổng hợp, dashboard, báo cáo và thông báo được triển khai sau để hoàn thiện trải nghiệm quản lý.

Trong từng nhóm chức năng, backend được ưu tiên định nghĩa domain rule và API trước, sau đó frontend tích hợp theo luồng người dùng. Những thay đổi schema được thực hiện qua migration mới. Những logic có rủi ro cao, như tính ngày nghỉ hoặc chuyển trạng thái đơn, được kiểm thử ở tầng service và API trước khi mở rộng giao diện.

\subsection{Chiến lược kiểm thử}
Chiến lược kiểm thử của hệ thống kết hợp nhiều mức để bao phủ cả logic nghiệp vụ và luồng sử dụng chính. Unit test tập trung vào các hàm và service có logic rõ ràng như tính ngày nghỉ, kiểm tra balance, validate trạng thái và xử lý quyền. Controller hoặc integration test kiểm tra API, validation, security và tương tác với cơ sở dữ liệu.

Frontend được kiểm tra bằng typecheck, lint, unit test cho schema/hook hoặc logic giao diện quan trọng, và build để phát hiện lỗi đóng gói. Với các luồng người dùng chính, smoke test bằng trình duyệt giúp xác nhận rằng người dùng có thể đăng nhập, nộp đơn, xử lý đơn và xem dữ liệu trên giao diện thật.

Bên cạnh test tự động, quá trình build LaTeX của báo cáo cũng được kiểm tra sau mỗi lần bổ sung nội dung lớn để phát hiện lỗi cú pháp, cảnh báo căn chỉnh hoặc sơ đồ tràn trang. Việc giữ source báo cáo ở dạng văn bản giúp các thay đổi nội dung, sơ đồ và bảng biểu có thể được theo dõi bằng Git như các phần khác của dự án.

\chapter{Kết quả thực hiện và đánh giá}

\section{Kết quả chức năng}
Sau quá trình triển khai, Leave Management System đã hình thành đầy đủ các nhóm chức năng chính của một hệ thống quản lý nghỉ phép trên nền web. Kết quả không chỉ dừng ở giao diện nhập liệu, mà bao gồm cả xử lý nghiệp vụ ở backend, lưu trữ dữ liệu bằng PostgreSQL, phân quyền theo vai trò, kiểm thử tự động và môi trường chạy bằng Docker. Các chức năng được xây dựng theo lát cắt hoàn chỉnh từ API đến giao diện để có thể demo theo luồng sử dụng thực tế.

\subsection{Xác thực, người dùng và dữ liệu nền}
Hệ thống đã hỗ trợ đăng nhập bằng email và mật khẩu, phát hành access token và refresh token, kiểm soát phiên làm việc và thu hồi refresh token khi cần. Mật khẩu được lưu dưới dạng hash, các API nghiệp vụ yêu cầu xác thực và quyền truy cập được kiểm tra theo vai trò. Cách tiếp cận này giúp frontend có thể hoạt động như một single page application trong khi backend vẫn giữ vai trò quyết định cuối cùng đối với quyền truy cập.

Nhóm chức năng quản trị dữ liệu nền bao gồm người dùng, phòng ban, loại nghỉ phép, ngày lễ và quỹ phép. Admin hoặc HR có thể quản lý các dữ liệu này để thiết lập tổ chức, xác định quản lý trực tiếp, cấu hình chính sách nghỉ phép và điều chỉnh quỹ phép theo nhu cầu nghiệp vụ. Đây là nền tảng để các luồng tạo đơn, duyệt đơn, tính ngày nghỉ và báo cáo hoạt động nhất quán.

\subsection{Quy trình đơn nghỉ phép}
Quy trình đơn nghỉ là phần nghiệp vụ trung tâm của hệ thống. Nhân viên có thể tạo đơn nghỉ với loại nghỉ, khoảng thời gian, lựa chọn nửa ngày và lý do. Backend tự tính số ngày nghỉ hợp lệ dựa trên ngày làm việc, cuối tuần và ngày lễ; đồng thời kiểm tra các điều kiện như ngày bắt đầu, ngày kết thúc, đơn chồng lấn và số ngày phép còn lại. Khi đơn hợp lệ, hệ thống lưu đơn ở trạng thái chờ duyệt và tạo thông tin cần thiết cho người có thẩm quyền xử lý.

Người quản lý, HR hoặc Admin có thể duyệt hoặc từ chối đơn theo phạm vi quyền hạn. Khi đơn được duyệt, hệ thống cập nhật trạng thái và trừ quỹ phép nếu loại nghỉ yêu cầu quản lý quota. Khi đơn bị từ chối hoặc bị hủy, hệ thống ghi nhận lý do hoặc hành động tương ứng, đồng thời đảm bảo quỹ phép được khôi phục trong các trường hợp cần thiết. Các thao tác quan trọng đều có lịch sử xử lý để phục vụ truy vết.

\subsection{Thông báo và tệp đính kèm}
Cơ chế thông báo trong ứng dụng đã được tích hợp vào vòng đời đơn nghỉ. Khi có sự kiện như nộp đơn, cập nhật đơn, phê duyệt, từ chối hoặc hủy đơn, hệ thống tạo thông báo cho người liên quan trong cùng transaction nghiệp vụ. Giao diện thông báo hỗ trợ phân biệt trạng thái đã đọc/chưa đọc, lọc thông báo, hiển thị loại sự kiện, thời điểm đọc và mở trực tiếp chi tiết đơn nghỉ liên quan.

Hệ thống cũng hỗ trợ tệp đính kèm cho đơn nghỉ trong phạm vi demo/test local. Metadata của file được lưu trong PostgreSQL, còn file vật lý được lưu trong Docker named volume. Quyền thao tác được giới hạn theo trạng thái và chủ thể của đơn: người tạo đơn có thể upload hoặc xóa file khi đơn còn chờ duyệt; người tạo, quản lý, HR và Admin có thể xem hoặc tải file nếu có quyền xem đơn. Chức năng này được chủ động giới hạn ở môi trường local để không kéo thêm dịch vụ lưu trữ ngoài vào phạm vi triển khai.

\subsection{Lịch, dashboard và báo cáo}
Lịch nghỉ phép đã hỗ trợ quan sát dữ liệu theo phạm vi người dùng. Employee và Manager có thể theo dõi lịch liên quan đến phòng ban hoặc nhóm của mình, trong khi HR/Admin có thêm khả năng xem toàn tổ chức hoặc lọc theo phòng ban và nhân viên. Mỗi mục lịch mang theo thông tin phòng ban để giao diện có thể tổng hợp và hiển thị rõ phạm vi đang xem.

Dashboard được cải thiện theo hướng cung cấp thông tin hành động thay vì chỉ liệt kê số liệu. Trang tổng quan có phần chào theo vai trò, các chỉ số cần chú ý, lối tắt đến thao tác thường dùng, biểu đồ quỹ phép và danh sách người đang nghỉ. Với HR/Admin, dashboard bổ sung góc nhìn tổ chức như số lượng nhân viên, phân bố trạng thái đơn và các phòng ban có nhiều đơn nghỉ trong kỳ.

Báo cáo đã được mở rộng từ xuất CSV cơ bản sang nhóm báo cáo có bộ lọc và preview. HR/Admin có thể xem nhanh tổng hợp số ngày nghỉ và số đơn theo tháng hoặc quý, loại nghỉ và phòng ban trước khi tải file. Các CSV chính gồm đơn nghỉ, quỹ phép và tổng hợp ngày nghỉ, giúp dữ liệu có thể được mở bằng công cụ bảng tính để phục vụ phân tích hoặc lưu trữ.

\begin{longtable}{p{0.24\textwidth}p{0.56\textwidth}p{0.12\textwidth}}
\toprule
\textbf{Nhóm kết quả} & \textbf{Nội dung đã hoàn thành} & \textbf{Trạng thái} \\
\midrule
\endfirsthead
\toprule
\textbf{Nhóm kết quả} & \textbf{Nội dung đã hoàn thành} & \textbf{Trạng thái} \\
\midrule
\endhead
Xác thực và phân quyền & Đăng nhập JWT, refresh token, RBAC, kiểm tra phạm vi dữ liệu theo vai trò và chủ thể nghiệp vụ. & Hoàn thành \\
Dữ liệu nền & Quản lý người dùng, phòng ban, loại nghỉ, ngày lễ, quỹ phép, quan hệ quản lý và dữ liệu seed demo. & Hoàn thành \\
Đơn nghỉ phép & Tạo, sửa, hủy, duyệt, từ chối, tính ngày nghỉ, kiểm tra chồng lấn, cập nhật quỹ phép và ghi lịch sử. & Hoàn thành \\
Thông báo & Notification in-app, badge chưa đọc, filter, mark-read và điều hướng tới chi tiết đơn. & Hoàn thành \\
Tệp đính kèm local-only & Upload, download, xóa file theo quyền; lưu metadata trong database và file trong Docker volume. & Hoàn thành \\
Lịch nghỉ & Lịch tháng/tuần, filter theo loại nghỉ, nhân viên, phòng ban và hiển thị phạm vi đang xem. & Hoàn thành \\
Dashboard & Tổng quan theo vai trò, quick actions, chỉ số cần chú ý, biểu đồ quỹ phép và số liệu tổ chức. & Hoàn thành \\
Báo cáo & CSV đơn nghỉ, CSV quỹ phép, CSV tổng hợp; preview theo tháng/quý, loại nghỉ và phòng ban. & Hoàn thành \\
\bottomrule
\end{longtable}

\section{Kết quả giao diện}
Frontend đã được xây dựng theo hướng ứng dụng nội bộ, ưu tiên tính rõ ràng, thao tác nhanh và thông tin nghiệp vụ dễ đọc. Các màn hình chính được tổ chức theo vai trò người dùng, giúp mỗi actor chỉ nhìn thấy các chức năng phù hợp với phạm vi công việc của mình. Giao diện sử dụng tiếng Việt, layout nhất quán, màu nhấn tiết chế và các trạng thái tải/rỗng/lỗi được xử lý để tránh cảm giác hệ thống bị treo khi dữ liệu chưa sẵn sàng.

\subsection{Luồng người dùng chính}
Ở phía nhân viên, các màn hình quan trọng gồm đăng nhập, hồ sơ cá nhân, nộp đơn nghỉ, danh sách đơn của tôi, chi tiết đơn, lịch nghỉ, dashboard và thông báo. Khi nộp đơn, giao diện hỗ trợ preview số ngày nghỉ, hiển thị gợi ý quỹ phép và cho phép đính kèm file trong môi trường local nếu chức năng được bật. Sau khi gửi đơn, người dùng có thể theo dõi trạng thái xử lý, xem lịch sử và nhận thông báo khi đơn thay đổi trạng thái.

Ở phía quản lý, giao diện tập trung vào danh sách đơn cần duyệt và thông tin đủ để ra quyết định: người nộp, loại nghỉ, thời gian nghỉ, số ngày tính toán, lý do và lịch sử xử lý. Các thao tác duyệt, từ chối hoặc hủy đơn được đặt trong ngữ cảnh chi tiết của đơn để giảm nhầm lẫn. Khi cần, quản lý có thể xem lịch nhóm hoặc dashboard để cân nhắc tình hình nhân sự trước khi xử lý.

Với HR/Admin, hệ thống cung cấp các màn hình quản trị dữ liệu nền, báo cáo, dashboard tổ chức và lịch toàn tổ chức. Các bộ lọc theo phòng ban, nhân viên, trạng thái, khoảng thời gian hoặc nhóm thời gian giúp người dùng khai thác dữ liệu mà không phải xuất toàn bộ rồi lọc thủ công bên ngoài.

\subsection{Cải thiện trải nghiệm sử dụng}
Một số cải thiện UI/UX đã được bổ sung để giao diện phù hợp hơn với demo và sử dụng thực tế. Dashboard mới có phần hero theo vai trò, quick actions, các card thể hiện việc cần chú ý, biểu đồ quỹ phép kèm thanh tiến độ và danh sách người đang nghỉ hôm nay rõ ràng hơn. Với HR/Admin, khối số liệu tổ chức giúp người dùng thấy nhanh tình trạng vận hành thay vì phải đi qua nhiều màn hình.

Các màn hình dùng dữ liệu từ API được bổ sung trạng thái tải, empty state và error state có nút thử lại. Điều này giúp người dùng phân biệt được trường hợp chưa có dữ liệu, đang tải dữ liệu hoặc request gặp lỗi. Một số thao tác nhạy cảm như khóa người dùng, hủy đơn hoặc xóa dữ liệu có bước xác nhận để hạn chế thao tác nhầm. Navigation trên thiết bị nhỏ cũng được cải thiện để demo trên nhiều kích thước màn hình thuận tiện hơn.

\begin{longtable}{p{0.24\textwidth}p{0.62\textwidth}}
\toprule
\textbf{Khu vực giao diện} & \textbf{Kết quả đạt được} \\
\midrule
\endfirsthead
\toprule
\textbf{Khu vực giao diện} & \textbf{Kết quả đạt được} \\
\midrule
\endhead
Dashboard & Hiển thị tổng quan theo vai trò, việc cần xử lý, quick actions, quỹ phép, người đang nghỉ và số liệu tổ chức cho HR/Admin. \\
Đơn nghỉ & Cho phép tạo, sửa, hủy, xem chi tiết, xem lịch sử xử lý và upload file đính kèm local-only khi điều kiện hợp lệ. \\
Phê duyệt & Cung cấp inbox đơn chờ xử lý, chi tiết đơn, thao tác duyệt/từ chối/hủy và phản hồi trạng thái rõ ràng. \\
Lịch & Hỗ trợ quan sát theo tháng/tuần, filter theo phòng ban hoặc nhân viên và tổng hợp số đơn theo phạm vi đang xem. \\
Báo cáo & Có form lọc, preview dữ liệu tổng hợp, thẻ chỉ số nhanh và nút tải CSV theo từng loại báo cáo. \\
Thông báo & Có chuông thông báo, badge chưa đọc, filter, trạng thái đọc/chưa đọc và điều hướng tới đơn liên quan. \\
Quản trị & Cung cấp màn hình quản lý người dùng, phòng ban, loại nghỉ, ngày lễ và quỹ phép với form validation. \\
\bottomrule
\end{longtable}

\section{Kết quả triển khai}
Hệ thống được đóng gói để có thể chạy toàn bộ stack bằng Docker Compose trong môi trường phát triển. Backend, frontend và PostgreSQL được khai báo trong cùng cấu hình compose, giúp quá trình khởi động, reset dữ liệu và kiểm thử thủ công nhất quán hơn giữa các máy. Frontend trong môi trường local được build và chạy bằng Vite preview để ổn định trên Docker Desktop Windows; backend chạy Spring Boot với cấu hình dev; database sử dụng PostgreSQL và migration được áp dụng bằng Flyway khi ứng dụng khởi động.

Bên cạnh môi trường local, dự án cũng có cấu hình phục vụ triển khai demo trực tuyến. Frontend được build thành static assets, backend được đóng gói thành container production, database sử dụng dịch vụ PostgreSQL quản lý, còn secrets được truyền qua biến môi trường. Các cấu hình như CORS, JWT secret, database URL, logging và security headers được tách khỏi mã nguồn để phù hợp với từng môi trường.

Một điểm cần tách bạch trong kết quả triển khai là chức năng tệp đính kèm. Do mục tiêu của chức năng này chỉ phục vụ demo/test local, hệ thống lưu file trong Docker named volume và giữ khả năng bật/tắt bằng cấu hình. Cách làm này đủ để minh họa nghiệp vụ đính kèm giấy tờ cho đơn nghỉ, nhưng không kéo thêm object storage hoặc dịch vụ hạ tầng bên ngoài vào phạm vi triển khai demo trực tuyến.

\begin{longtable}{p{0.25\textwidth}p{0.61\textwidth}}
\toprule
\textbf{Hạng mục} & \textbf{Kết quả triển khai} \\
\midrule
\endfirsthead
\toprule
\textbf{Hạng mục} & \textbf{Kết quả triển khai} \\
\midrule
\endhead
Container hóa & Backend, frontend và PostgreSQL chạy được bằng Docker Compose; cấu hình dev và production được tách riêng. \\
Database migration & Flyway quản lý lịch sử schema, giúp khởi tạo database lặp lại được và tránh sửa trực tiếp schema đã áp dụng. \\
Cấu hình môi trường & Các giá trị nhạy cảm như database URL, JWT secret, CORS origin và SMTP được đưa ra biến môi trường. \\
Demo data & Dữ liệu seed phục vụ demo gồm người dùng, phòng ban, quỹ phép và đơn nghỉ ở nhiều trạng thái. \\
Local-only attachment & Metadata lưu trong PostgreSQL, file lưu trong Docker volume; có cờ cấu hình để giới hạn phạm vi sử dụng. \\
\bottomrule
\end{longtable}

\section{Kết quả kiểm thử}
Kiểm thử được thực hiện ở nhiều mức để bao phủ cả logic nghiệp vụ và giao diện. Với backend, các test tập trung vào service, controller và luồng tích hợp quan trọng như đăng nhập, quản lý danh mục, tính ngày nghỉ, tạo đơn, cập nhật đơn, phê duyệt, từ chối, hủy đơn, cập nhật quỹ phép, thông báo, báo cáo và quyền truy cập. Các tình huống có rủi ro cao như đơn chồng lấn, ngày nghỉ quá khứ, thiếu quỹ phép hoặc thao tác sai vai trò được kiểm tra ở tầng backend để đảm bảo dữ liệu không bị phụ thuộc vào validation phía frontend.

Với frontend, dự án sử dụng typecheck, lint, unit test và build để phát hiện lỗi kiểu dữ liệu, lỗi quy ước code và lỗi bundle trước khi chạy demo. Các test giao diện tập trung vào utility, schema form, component quan trọng và các trạng thái tải/lỗi/rỗng. Ngoài ra, smoke test bằng trình duyệt được dùng để kiểm tra các luồng chính theo vai trò như đăng nhập, nộp đơn, xử lý đơn, xem lịch, mở thông báo và truy cập dashboard.

Trong lần kiểm tra gần nhất của báo cáo, các quality gate chính đã được chạy lại: frontend typecheck, lint, unit test, build; backend/API được kiểm tra qua container đang chạy; LaTeX report được build thành công; và các request health/frontend local trả về trạng thái hợp lệ. Các cảnh báo lint còn lại thuộc nhóm cảnh báo kỹ thuật đã biết của một số component dùng chung, không phải lỗi build hoặc lỗi chức năng.

\begin{longtable}{p{0.22\textwidth}p{0.48\textwidth}p{0.18\textwidth}}
\toprule
\textbf{Loại kiểm tra} & \textbf{Phạm vi} & \textbf{Kết quả} \\
\midrule
\endfirsthead
\toprule
\textbf{Loại kiểm tra} & \textbf{Phạm vi} & \textbf{Kết quả} \\
\midrule
\endhead
Backend test & Logic nghiệp vụ, API, bảo mật, dữ liệu và các luồng xử lý đơn nghỉ. & Đạt \\
Frontend typecheck & Kiểm tra TypeScript cho component, hook, API client và schema. & Đạt \\
Frontend lint & Kiểm tra quy ước code frontend; còn một số cảnh báo đã biết ở component dùng chung. & Đạt \\
Frontend unit test & Utility, schema form, component và trạng thái giao diện quan trọng. & Đạt \\
Frontend build & Kiểm tra khả năng build production bundle bằng Vite. & Đạt \\
Smoke test & Kiểm tra luồng người dùng chính trên trình duyệt/local stack. & Đạt \\
LaTeX build & Kiểm tra cú pháp, bảng biểu và khả năng xuất PDF của báo cáo. & Đạt \\
\bottomrule
\end{longtable}

\section{Đánh giá}
Nhìn chung, hệ thống đã đáp ứng được mục tiêu chính của mini-project: mô phỏng một quy trình quản lý nghỉ phép hoàn chỉnh bằng ứng dụng web có frontend, backend, database, phân quyền, kiểm thử và cấu hình triển khai. Các chức năng cốt lõi đã đi qua đủ vòng đời từ phân tích yêu cầu đến hiện thực và kiểm tra, trong đó các phần quan trọng như tính ngày nghỉ, cập nhật quỹ phép, chuyển trạng thái đơn và phân quyền được đặt ở backend thay vì chỉ xử lý ở giao diện.

\subsection{Mức độ đáp ứng yêu cầu}
Các yêu cầu nghiệp vụ chính đã được hiện thực ở mức đủ để demo theo luồng thực tế. Nhân viên có thể tự phục vụ các thao tác liên quan đến đơn nghỉ; quản lý có công cụ xử lý đơn và theo dõi nhóm; HR/Admin có khả năng quản lý dữ liệu nền, xem lịch toàn tổ chức và khai thác báo cáo. Hệ thống cũng có các cơ chế hỗ trợ vận hành như audit, notification, dashboard, seed dữ liệu demo và các bộ lọc dữ liệu.

So với mục tiêu ban đầu là xây dựng một sản phẩm học tập nhưng có cấu trúc kỹ thuật nghiêm túc, dự án đã thể hiện được các phần quan trọng của phát triển phần mềm web: thiết kế API, tổ chức module theo feature, quản lý schema bằng migration, kiểm soát security, quản lý server state ở frontend và kiểm thử tự động. Các cải thiện gần đây như báo cáo nâng cao, lịch theo phòng ban, notification tốt hơn và dashboard rõ ràng hơn giúp sản phẩm thuyết phục hơn khi trình bày trong bối cảnh thực tập hoặc portfolio.

\subsection{Ưu điểm}
Ưu điểm lớn nhất của hệ thống là ranh giới trách nhiệm rõ ràng. Backend chịu trách nhiệm tính toán và bảo vệ nghiệp vụ; frontend tập trung vào trải nghiệm thao tác; PostgreSQL lưu trữ dữ liệu có ràng buộc; Flyway quản lý thay đổi schema; Docker Compose giúp môi trường chạy nhất quán. Cấu trúc package theo feature cũng giúp việc thêm chức năng mới như notification, report, calendar hoặc attachment không làm rối toàn bộ codebase.

Một ưu điểm khác là hệ thống có tính demo tốt. Dữ liệu seed tương đối tự nhiên, giao diện đã có các luồng theo vai trò, dashboard và báo cáo giúp người xem hiểu nhanh giá trị nghiệp vụ. Các tính năng local-only như attachment được giới hạn rõ phạm vi, đủ để minh họa ý tưởng mà không làm tăng chi phí hoặc độ phức tạp hạ tầng.

\subsection{Hạn chế}
Do đây là mini-project, hệ thống chưa hướng tới đầy đủ các yêu cầu vận hành của một sản phẩm doanh nghiệp quy mô lớn. Chức năng tệp đính kèm mới dừng ở local storage, chưa dùng object storage, antivirus scanning hoặc CDN. Notification hiện là in-app và email local/dev, chưa có kênh thời gian thực như WebSocket hoặc push notification. Báo cáo sử dụng CSV và preview tổng hợp, chưa có dashboard phân tích sâu hoặc xuất Excel nhiều sheet.

Một số khía cạnh vận hành như quan sát hệ thống, phân tích log dài hạn, backup tự động, phân quyền rất chi tiết theo từng permission hoặc tích hợp SSO chưa được đưa vào phạm vi hiện tại. Đây là lựa chọn có chủ đích để giữ dự án tập trung vào bài toán quản lý nghỉ phép và tránh mở rộng quá mức so với mục tiêu mini-project.

\subsection{Nhận xét tổng quan}
Kết quả hiện tại cho thấy dự án đã đạt trạng thái đủ hoàn chỉnh để trình bày như một bài toán software development/web programming: có bối cảnh nghiệp vụ, yêu cầu, mô hình dữ liệu, luồng xử lý, kiến trúc triển khai, giao diện theo vai trò, kiểm thử và đánh giá giới hạn. Phần còn lại của báo cáo nên tiếp tục tập trung vào việc diễn giải nhất quán các quyết định kỹ thuật, minh họa kết quả bằng hình ảnh giao diện chọn lọc và tổng kết bài học thay vì bổ sung thêm nhiều chức năng mới.

\chapter{Kết luận}

\section{Kết quả đạt được}
Tóm tắt những kết quả chính đạt được sau quá trình xây dựng Leave Management System.

\section{Bài học kỹ thuật}
TODO: Tổng kết các bài học về phát triển ứng dụng web, thiết kế backend/frontend, quản lý database migration, kiểm thử và triển khai.

\section{Hạn chế}
Nêu các hạn chế đã xác định, bao gồm phạm vi demo, chi phí hạ tầng free tier và các tính năng tương lai chỉ phục vụ kiểm thử local.

\section{Hướng phát triển}
Đề xuất hướng phát triển tiếp theo, nếu tiếp tục demo local, chủ yếu là tinh chỉnh giao diện theo phản hồi sử dụng thực tế và duy trì các chức năng đã có. Hướng này không mở rộng sang các bài toán nghiệp vụ phức tạp hơn phạm vi hiện tại.

\begin{thebibliography}{99}
\phantomsection
\sloppy
\addcontentsline{toc}{chapter}{Tài liệu tham khảo}

\bibitem{spring-boot}
VMware, ``Spring Boot Reference Documentation'', \url{https://docs.spring.io/spring-boot/}.

\bibitem{spring-security}
VMware, ``Spring Security Reference Documentation'', \url{https://docs.spring.io/spring-security/reference/}.

\bibitem{react}
Meta Open Source, ``React Documentation'', \url{https://react.dev/}.

\bibitem{postgresql}
PostgreSQL Global Development Group, ``PostgreSQL Documentation'', \url{https://www.postgresql.org/docs/}.

\bibitem{docker}
Docker Inc., ``Docker Documentation'', \url{https://docs.docker.com/}.

\end{thebibliography}


\appendix
\chapter{Phụ lục}

\section{Nguồn tài liệu nội bộ}
Các tài liệu trong kho mã nguồn được dùng làm nguồn chính khi hoàn thiện báo cáo:

\begin{itemize}
  \item \texttt{README.md}: mô tả tổng quan dự án, cách chạy, tài khoản demo và ảnh minh họa.
  \item \texttt{docs/REQUIREMENTS.md}: mô tả yêu cầu nghiệp vụ, phạm vi và các quyết định sản phẩm.
  \item \texttt{docs/ARCHITECTURE.md}: mô tả kiến trúc backend/frontend/database, bảo mật, notification, attachment và báo cáo.
  \item \texttt{docs/DATABASE.md}: mô tả schema logic, quan hệ bảng và migration.
  \item \texttt{docs/DEPLOYMENT.md}: mô tả cấu hình triển khai demo và các lưu ý vận hành.
  \item \texttt{docs/ROADMAP.md}: mô tả tiến độ các nhóm chức năng và định hướng demo/test local.
  \item \texttt{docs/UI-GUIDELINES.md}: mô tả định hướng giao diện tiếng Việt, màu sắc và quy tắc trình bày.
\end{itemize}

\section{Tóm tắt API chính}
Bảng dưới đây liệt kê các nhóm endpoint chính ở mức tra cứu. Chi tiết request/response được giữ trong mã nguồn backend và các DTO tương ứng.

\begingroup
\small
\begin{longtable}{p{0.20\textwidth}p{0.30\textwidth}p{0.36\textwidth}}
\toprule
\textbf{Nhóm API} & \textbf{Endpoint tiêu biểu} & \textbf{Mục đích} \\
\midrule
\endfirsthead
\toprule
\textbf{Nhóm API} & \textbf{Endpoint tiêu biểu} & \textbf{Mục đích} \\
\midrule
\endhead
Auth & \path{/auth/*} & Đăng nhập, refresh token, đăng xuất và lấy thông tin người dùng hiện tại. \\
User & \path{/users/*} & Quản lý tài khoản, hồ sơ cá nhân, trạng thái hoạt động và đổi/đặt lại mật khẩu. \\
Department & \path{/departments/*} & Quản lý phòng ban và danh bạ phòng ban của người dùng. \\
Leave type, holiday & \path{/leave-types/*}; \path{/holidays/*} & Quản lý loại nghỉ phép và ngày lễ phục vụ tính ngày nghỉ. \\
Leave balance & \path{/leave-balances/*} & Xem, khởi tạo, điều chỉnh và carry-over quỹ phép. \\
Leave request & \path{/leave-requests/*} & Tạo, sửa, xem danh sách, xem chi tiết và lịch sử đơn nghỉ. \\
Approval & \path{/leave-requests/*/approve}; \path{/*/reject}; \path{/*/cancel} & Duyệt, từ chối hoặc hủy đơn nghỉ theo quyền và trạng thái hợp lệ. \\
Calendar, dashboard & \path{/calendar}; \path{/dashboard/*} & Cung cấp dữ liệu lịch nghỉ và số liệu tổng quan theo vai trò. \\
Report & \path{/reports/*} & Preview và tải CSV cho báo cáo đơn nghỉ, tổng hợp ngày nghỉ và quỹ phép. \\
Notification & \path{/notifications/*} & Xem, đếm và đánh dấu thông báo đã đọc. \\
Attachment local-only & \path{/leave-requests/*/attachments/*} & Upload, liệt kê, tải và xóa tệp đính kèm trong môi trường local. \\
\bottomrule
\end{longtable}
\endgroup

\section{Schema cơ sở dữ liệu rút gọn}
Bảng dưới đây tóm tắt các bảng dữ liệu chính, khóa ngoại quan trọng và vai trò lưu trữ. Các cột audit dùng chung không được lặp lại ở mọi dòng.

\begin{longtable}{p{0.22\textwidth}p{0.30\textwidth}p{0.34\textwidth}}
\toprule
\textbf{Bảng} & \textbf{Khóa/quan hệ chính} & \textbf{Dữ liệu quan trọng} \\
\midrule
\endfirsthead
\toprule
\textbf{Bảng} & \textbf{Khóa/quan hệ chính} & \textbf{Dữ liệu quan trọng} \\
\midrule
\endhead
\texttt{departments} & PK \texttt{id}; \texttt{head\_user\_id -> users} & Mã phòng ban, tên phòng ban, trạng thái hoạt động. \\
\texttt{users} & PK \texttt{id}; \texttt{department\_id -> departments}; \texttt{manager\_id -> users} & Email, mã nhân viên, hash mật khẩu, họ tên, vai trò, ngày vào làm. \\
\texttt{leave\_types} & PK \texttt{id}; unique \texttt{code} & Mã loại nghỉ, quota mặc định, cờ yêu cầu quỹ phép, trạng thái hoạt động. \\
\texttt{leave\_balances} & FK tới \texttt{users}, \texttt{leave\_types}; unique theo user/type/year & Năm, tổng ngày, đã dùng, điều chỉnh, chuyển tiếp, số ngày còn lại tính ở app. \\
\texttt{holidays} & Unique \texttt{holiday\_date} & Ngày lễ và mô tả, dùng khi tính số ngày làm việc. \\
\texttt{leave\_requests} & FK tới \texttt{users}, \texttt{leave\_types}, \texttt{manager\_id -> users} & Khoảng ngày, nửa ngày, tổng ngày, lý do, trạng thái, quản lý tại thời điểm tạo. \\
\texttt{approval\_actions} & FK tới \texttt{leave\_requests}, \texttt{actor\_id -> users} & Hành động, trạng thái trước/sau, ghi chú xử lý. \\
\texttt{notifications} & FK tới \texttt{users}, \texttt{leave\_requests} & Loại sự kiện, nội dung thông báo, trạng thái đọc, thời điểm đọc. \\
\texttt{attachments} & FK tới \texttt{leave\_requests}, \texttt{uploaded\_by -> users}; unique \texttt{stored\_key} & Tên file gốc, key lưu trữ, content type, dung lượng; file vật lý nằm trong Docker volume. \\
\texttt{audit\_log} & FK tùy chọn tới \texttt{users} & Actor, hành động, loại đối tượng, id đối tượng, dữ liệu trước/sau. \\
\texttt{refresh\_tokens} & FK tới \texttt{users}; unique \texttt{token\_hash} & Hash refresh token, thời điểm hết hạn và thời điểm thu hồi. \\
\bottomrule
\end{longtable}

\section{Hình ảnh giao diện bổ sung}
Phần thân báo cáo đã sử dụng các ảnh chọn lọc cho dashboard, đơn nghỉ, phê duyệt, notification, lịch và báo cáo. Hình dưới đây bổ sung minh họa cho giao diện quản trị người dùng, một phần quan trọng trong nhóm chức năng Admin/HR.

\begin{figure}[H]
\centering
\includegraphics[width=0.92\textwidth]{images/users.png}
\caption{Giao diện quản trị người dùng}
\end{figure}


\end{document}
